\begin{resumo}[Abstract]
 \begin{otherlanguage*}{english}
  The organization of amateur football matches in Brasília faces recurring challenges, particularly in the public courts spread across the city. Games are usually arranged through dispersed messaging groups, resulting in inconsistent confirmations and little clarity about space availability, which leads to cancellations, missed connections, and underuse of public facilities. In response, this work proposes Borafut, a cross-platform mobile application that streamlines the organization of amateur football in the Federal District by connecting players, courts, and local groups through intuitive and accessible features. The development was guided by a hybrid methodological approach that combines plan-driven elements and agile practices, grounded in Extreme Programming (XP). The requirements gathering process relied on observation, brainwriting, and a questionnaire answered by more than 60 amateur players, while prioritization and modeling adopted the MoSCoW method and a hierarchical backlog decomposition spanning successive levels of granularity, from strategic themes down to technical tasks. The aim of this work is to cover the complete software development cycle for a minimum viable product, from problem investigation and requirements definition to the design of the system architecture, database, and visual identity. This effort produced a high-fidelity interactive prototype, used as a proof of concept for Borafut's main user journeys, and culminated in the implementation of the minimum viable product: a cross-platform mobile application and a backend service with a relational database, real-time communication, and an automated test suite subject to a coverage gate, submitted to validation with real users.

  \vspace{\onelineskip}

  \noindent 
  \textbf{Keywords}: amateur football. mobile application. software engineering. agile methodologies. public spaces.
  \end{otherlanguage*}
\end{resumo}